% !LW recipe = make-deck
% Level 5T (proofs track, TPIL): induction on lists. Theory: list programs
% as equations; induction; a proof. Lean: the programs (given), the first
% induction (worked), then exercises.
% Build: make (rules in ../common/deck.mk).
\DocumentMetadata{
  lang        = en-GB,
  pdfstandard = ua-2,
  tagging     = on,
}
\documentclass{ltx-talk}

\input{preamble}

\title{Level 5T: Induction on Lists}
\author{The Proofs and Programs Club}
\institute{Royal Holloway, University of London}

\begin{document}

\leveltitle{5T}{Induction on Lists}{Proofs track: follows TPIL}
  {After \levellink{4t}}{Next: \levellink{6}, with \levellink{5f}}

\levelframe{5T}{Level 5T at a Glance}{%
  nothing new: the list programs \code{select}, \code{project} and
  \code{product} are given (Levels 2F and 3F build them).}{%
  facts about them for every list, by induction.}{%
  prove by induction, step by step & \rds{nng} (Addition World) \\
  rewrite and simplify: \code{rw}, \code{simp} & \rds{tpil-rw, tpil-simp}
    \\
  prove by induction on lists, and split on an \code{if} &
    \rds{tpil-indtac, tpil-split} \\
  prove facts about multiplication & \rds{nng} (Multiplication World) \\
}

% ---------------------------------------------------------------- theory
\theorypart{List programs, and induction}

\begin{frame}
  \frametitle{A List Program Is a Few Equations}
  Write $t :: l$ for the list with first row $t$, then the list $l$.
  \defn{Definition 5T.1 (Selection and projection on lists)}{
    \[
      \begin{array}{@{}l@{\;=\;}l@{\qquad}l@{}}
        \mathit{select}\ p\ [\,] & [\,] \\
        \mathit{select}\ p\ (t :: l) & t :: \mathit{select}\ p\ l &
          \text{if } p(t) = \true \\
        \mathit{select}\ p\ (t :: l) & \mathit{select}\ p\ l &
          \text{if } p(t) = \false \\
        \mathit{project}\ f\ [\,] & [\,] \\
        \mathit{project}\ f\ (t :: l) & f(t) :: \mathit{project}\ f\ l \\
      \end{array}
    \]}
  Here $p$ is a condition with values in $\B = \{\true, \false\}$, as a
  program needs. These lists represent $\sel{p}(R)$ and $f[R]$ (Level 6
  proves it).
  \theorytag{5T}
\end{frame}

\begin{frame}
  \frametitle{Induction Proves a Fact About Every List}
  \defn{Principle (Induction on lists)}{To prove $P(l)$ for every list
    $l$: prove $P([\,])$; and prove $P(t :: l)$ for every row $t$ and
    list $l$, assuming $P(l)$.}
  It is the induction of the Natural Number Game: $[\,]$ plays $0$, and
  $t :: l$ plays $n + 1$.
  \defn{Fact 5T.2}{$\mathit{length}(\mathit{project}\ f\ l) =
    \mathit{length}(l)$, for every list $l$.}
  \emph{Proof.} For $[\,]$, both sides are $0$. For $t :: l$: the length
  of $\mathit{project}\ f\ (t :: l) = f(t) :: \mathit{project}\ f\ l$ is
  $1 + \mathit{length}(\mathit{project}\ f\ l)$, which is
  $1 + \mathit{length}(l)$ by assumption. \proofend
  \theorytag{5T}
\end{frame}

\begin{frame}
  \frametitle{Law 1 Holds for Lists, Too}
  \defn{Fact 5T.3}{$\mathit{project}\ f\ (\mathit{project}\ g\ l) =
    \mathit{project}\ (f \circ g)\ l$.}
  \defn{Fact 5T.4 (Law 1, for lists)}{$\mathit{select}\ p\
    (\mathit{select}\ q\ l) = \mathit{select}\ (p \wedge q)\ l$, as lists:
    the same rows, in the same order.}
  \emph{Proof.} By induction on $l$. For $[\,]$, both sides are $[\,]$.
  For $t :: l$, take the cases of $q(t)$ and $p(t)$:
  \begin{itemize}
    \item if $q(t) = \false$, both sides drop $t$, and the assumption
      finishes;
    \item if $q(t) = \true$, both sides keep $t$ exactly when $p(t) =
      \true$. \proofend
  \end{itemize}
  \theorytag{5T}
\end{frame}

% ---------------------------------------------------------------- Lean
\leanpart{Your first proofs by induction}

\begin{frame*}
  \frametitle{The Programs, in Lean}
  \leanblock{lists.lean}
  \begin{itemize}
    \item Definition 5T.1, and Level 3F's product: copy them, or use
      yours.
  \end{itemize}
  \slideref{\lnk{tpil-match}}
  \leantag{5T}
\end{frame*}

\begin{frame*}
  \frametitle{Fact 5T.2 in Lean: Your First Proof by Induction}
  \leanblock{induction.lean}
  \begin{itemize}
    \item The same proof as on paper: one case for \code{nil}, one for
      \code{cons}. \code{simp} unfolds \code{project}; \code{ih} is the
      assumption $P(l)$.
    \item \ex{5T.1} Replace \code{simp [project, ih]} with \code{rw}
      steps, and watch the goal change in the Infoview.
  \end{itemize}
  \slideref{\lnk{tpil-indtac} \rfsep \lnk{nng} (Addition World)}
  \leantag{5T}
\end{frame*}

\begin{frame*}
  \frametitle{Your Turn: Projections Compose}
  \leanblock{compose.lean}
  \begin{itemize}
    \item \ex{5T.2} Fact 5T.3: the same shape as the proof before.
      Induction on \code{l}, then \code{simp [project, ih]}.
  \end{itemize}
  \slideref{\lnk{tpil-indtac} \rfsep \lnk{tpil-simp}}
  \leantag{5T}
\end{frame*}

\begin{frame*}
  \frametitle{Your Turn: Law 1 for Lists}
  \leanblock{law1.lean}
  \begin{itemize}
    \item \ex{5T.3} Fact 5T.4. Induction on \code{l}. In the \code{cons}
      case, take the cases of the conditions, as on paper: \code{cases hq
      : q row} and \code{cases hp : p row}.
    \item Then \code{simp [select, hp, hq, ih]} finishes each case.
  \end{itemize}
  \slideref{\lnk{tpil-indtac} \rfsep \lnk{tpil-split}}
  \leantag{5T}
\end{frame*}

\begin{frame*}
  \frametitle{Your Turn: A Product of m Rows and n Rows Has m × n}
  \leanblock{size.lean}
  \begin{itemize}
    \item \ex{5T.4} Induction on \code{r}, then \code{simp [product,
      project\_length, ih]}, using Fact 5T.2.
    \item What is left is arithmetic: \code{Nat.add\_mul},
      \code{Nat.one\_mul} and \code{Nat.add\_comm} finish it.
  \end{itemize}
  \slideref{\lnk{nng} (Multiplication World) \rfsep \lnk{tpil-rw}}
  \leantag{5T}
\end{frame*}

\end{document}
